\appendix

\section{Reproducibility Appendix}
\label{sec:appendix}

\subsection{Software Version Matrix}
\label{sec:appendix:versions}

\begin{table}[h]
\centering
\caption{Pinned software versions for all experiments.  ROCm and CUDA wheels
are not interchangeable; the \texttt{torch} version strings differ only in the
\texttt{+cu12x} / \texttt{+rocm7.1} suffix.}
\label{tab:software_versions}
\small
\begin{tabular}{lll}
\toprule
\textbf{Component} & \textbf{NVIDIA (H100)} & \textbf{AMD (MI350)} \\
\midrule
OS                 & Ubuntu 22.04           & Ubuntu 22.04 \\
Kernel             & \PH{}                  & \PH{} \\
Driver             & \PH{} (CUDA 12.x)      & ROCm \PH{} \\
\texttt{torch}     & \texttt{2.10.0+cu12x}  & \texttt{2.10.0+rocm7.1} \\
\texttt{torchvision} & \PH{}                & \PH{} \\
\texttt{transformers} & \PH{}               & \PH{} \\
\texttt{tokenizers} & \PH{}                 & \PH{} \\
\texttt{vllm}      & \PH{} (commit \PH{})   & \PH{} (commit \PH{}) \\
\texttt{sglang}    & \PH{} (commit \PH{})   & \PH{} (commit \PH{}) \\
\texttt{aiohttp}   & \PH{}                  & \PH{} \\
\texttt{swanlab}   & \PH{}                  & \PH{} \\
\texttt{triton}    & \PH{}                  & \PH{} (ROCm Triton) \\
MiniMind commit    & \PH{}                  & same \\
\bottomrule
\end{tabular}
\end{table}

\subsection{Environment Setup}
\label{sec:appendix:setup}

\paragraph{NVIDIA (H100).}

\begin{verbatim}
# Install PyTorch 2.10 CUDA 12.x wheel
pip install torch==2.10.0+cu12x torchvision \
    --index-url https://download.pytorch.org/whl/cu12x

# Install remaining dependencies
pip install transformers tokenizers aiohttp swanlab \
            fastapi uvicorn requests

# Install inference engines (CUDA)
pip install vllm==<PIN>
pip install sglang[all]==<PIN>
\end{verbatim}

\paragraph{AMD (MI350).}

\begin{verbatim}
# ROCm 7.1 must be pre-installed via AMD's repo
# Verify: rocm-smi --version

# Install PyTorch 2.10 ROCm wheel
pip install torch==2.10.0+rocm7.1 torchvision \
    --index-url https://download.pytorch.org/whl/rocm7.1

# Same Python dependencies
pip install transformers tokenizers aiohttp swanlab \
            fastapi uvicorn requests

# Install inference engines (ROCm)
# vLLM ROCm: build from source or use AMD-provided wheel
pip install vllm==<PIN_ROCM>
pip install sglang[all]==<PIN_ROCM>

# Optional: override GFX version for RDNA3 dev machines
# (not needed on MI350 gfx942)
# export HSA_OVERRIDE_GFX_VERSION=11.0.0
\end{verbatim}

\subsection{Dataset Preparation}
\label{sec:appendix:datasets}

All datasets are gitignored and must be downloaded separately.  Checksums
below allow verification that the correct files are in place before running
any experiment block.

\begin{table}[h]
\centering
\caption{Dataset files, sizes, and SHA-256 hashes (first 8 hex digits).}
\label{tab:datasets}
\small
\begin{tabular}{llrr}
\toprule
\textbf{File} & \textbf{Stage} & \textbf{Size} & \textbf{SHA-256 (first 8)} \\
\midrule
\texttt{dataset/pretrain\_t2t\_mini.jsonl} & Pretrain  & \PH{} & \PH{} \\
\texttt{dataset/sft\_t2t\_mini.jsonl}      & SFT       & \PH{} & \PH{} \\
\texttt{dataset/dpo.jsonl}                 & DPO       & \PH{} & \PH{} \\
\texttt{dataset/lora\_medical.jsonl}       & LoRA      & \PH{} & \PH{} \\
\texttt{dataset/rlaif.jsonl}               & PPO/GRPO  & \PH{} & \PH{} \\
\texttt{dataset/agent\_rl.jsonl}           & Agent RL  & \PH{} & \PH{} \\
\texttt{out/internlm2-1\_8b-reward/}       & PPO/GRPO reward & \PH{} & \PH{} \\
\bottomrule
\end{tabular}
\end{table}

\subsection{Canonical Run Commands}
\label{sec:appendix:commands}

All trainer commands are run from the \texttt{trainer/} directory.  All
evaluation and conversion commands are run from the repo root (or the
\texttt{scripts/} directory as noted).

\subsubsection{Block A: Pretrain}

\begin{verbatim}
# Dense, 1 GPU, bf16, no compile
torchrun --nproc_per_node=1 train_pretrain.py \
    --data_path ../dataset/pretrain_t2t_mini.jsonl \
    --hidden_size 768 --num_hidden_layers 8 \
    --epochs 2 --batch_size 32 --dtype bfloat16 \
    --use_compile 0 --use_wandb 0 \
    --save_weight pretrain 2>&1 | tee ../paper/logs/block_a/<VENDOR>_dense_bf16_g1_nocompile.log

# Dense, 8 GPUs, bf16, compile
torchrun --nproc_per_node=8 train_pretrain.py \
    --data_path ../dataset/pretrain_t2t_mini.jsonl \
    --hidden_size 768 --num_hidden_layers 8 \
    --epochs 2 --batch_size 32 --dtype bfloat16 \
    --use_compile 1 --use_wandb 0 \
    --save_weight pretrain 2>&1 | tee ../paper/logs/block_a/<VENDOR>_dense_bf16_g8_compile.log

# MoE, 8 GPUs, bf16
torchrun --nproc_per_node=8 train_pretrain.py \
    --data_path ../dataset/pretrain_t2t_mini.jsonl \
    --hidden_size 768 --num_hidden_layers 8 \
    --use_moe 1 --epochs 2 --batch_size 32 \
    --dtype bfloat16 --use_compile 0 --use_wandb 0 \
    --save_weight pretrain_moe 2>&1 | tee ../paper/logs/block_a/<VENDOR>_moe_bf16_g8.log
\end{verbatim}

\subsubsection{Block B: Post-training}

\begin{verbatim}
# SFT, 8 GPUs
torchrun --nproc_per_node=8 train_full_sft.py \
    --from_weight pretrain \
    --data_path ../dataset/sft_t2t_mini.jsonl \
    --hidden_size 768 --num_hidden_layers 8 \
    --epochs 1 --batch_size 16 --dtype bfloat16 \
    --save_weight full_sft 2>&1 | tee ../paper/logs/block_b/<VENDOR>_sft_g8.log

# LoRA, 1 GPU
torchrun --nproc_per_node=1 train_lora.py \
    --from_weight full_sft \
    --data_path ../dataset/lora_medical.jsonl \
    --hidden_size 768 --lora_name lora_medical \
    --dtype bfloat16 2>&1 | tee ../paper/logs/block_b/<VENDOR>_lora_g1.log

# Distillation: MoE teacher -> dense student
torchrun --nproc_per_node=8 train_distillation.py \
    --teacher_weight pretrain_moe --student_weight full_sft \
    --hidden_size 768 --use_moe 1 --dtype bfloat16 \
    --alpha 0.5 --temperature 1.5 \
    2>&1 | tee ../paper/logs/block_b/<VENDOR>_distill_g8.log
\end{verbatim}

\subsubsection{Block C: Preference Alignment}

\begin{verbatim}
# DPO
torchrun --nproc_per_node=1 train_dpo.py \
    --from_weight full_sft \
    --data_path ../dataset/dpo.jsonl \
    --hidden_size 768 --dtype bfloat16 \
    --beta 0.15 --learning_rate 4e-8 \
    2>&1 | tee ../paper/logs/block_c/<VENDOR>_dpo_g1.log

# GRPO, torch rollout engine, num_generations=4
torchrun --nproc_per_node=8 train_grpo.py \
    --from_weight full_sft \
    --data_path ../dataset/rlaif.jsonl \
    --hidden_size 768 --dtype bfloat16 \
    --num_generations 4 --rollout_engine torch \
    --reward_model_path ../out/internlm2-1_8b-reward \
    2>&1 | tee ../paper/logs/block_c/<VENDOR>_grpo_torch_g4.log

# PPO, sglang rollout engine
torchrun --nproc_per_node=8 train_ppo.py \
    --from_weight full_sft \
    --data_path ../dataset/rlaif.jsonl \
    --hidden_size 768 --dtype bfloat16 \
    --rollout_engine sglang \
    --reward_model_path ../out/internlm2-1_8b-reward \
    2>&1 | tee ../paper/logs/block_c/<VENDOR>_ppo_sglang_g8.log
\end{verbatim}

\subsubsection{Block D: Agent RL}

\begin{verbatim}
torchrun --nproc_per_node=8 train_agent.py \
    --from_weight full_sft \
    --data_path ../dataset/agent_rl.jsonl \
    --hidden_size 768 --dtype bfloat16 \
    --reward_model_path ../out/internlm2-1_8b-reward \
    2>&1 | tee ../paper/logs/block_d/<VENDOR>_agent_g8.log
\end{verbatim}

\subsubsection{Block E: Inference}

\begin{verbatim}
# Convert to Qwen3 HF format (run from repo root)
python scripts/convert_model.py \
    --weight full_sft --hidden_size 768 \
    --out_path out/qwen3_converted/full_sft_blockb_dense_g8

# HF-native single-stream benchmark
python eval_llm.py --weight full_sft --hidden_size 768 \
    --show_speed 1 --max_new_tokens 512

# Concurrent benchmark via OpenAI-compatible server
cd scripts && python serve_openai_api.py \
    --weight full_sft --hidden_size 768 --port 8998 &
sleep 5
python ../paper/experiments/bench_client.py \
    --port 8998 --concurrency 32 --num_requests 200 --max_tokens 512

# vLLM serve (Qwen3-converted weights)
python -m vllm.entrypoints.openai.api_server \
    --model out/qwen3_converted/full_sft_blockb_dense_g8 \
    --port 8999 --tensor-parallel-size 1 &
sleep 10
python paper/experiments/bench_client.py \
    --port 8999 --concurrency 32 --num_requests 200 --max_tokens 512

# Verification: greedy parity (Gate 6)
python paper/experiments/verify_greedy_parity.py \
    --native_weight full_sft_blockb_dense_g8 \
    --conv_path out/qwen3_converted/full_sft_blockb_dense_g8 \
    --server_url http://127.0.0.1:8999 --num_prompts 20

# Tool-call eval
python scripts/eval_toolcall.py \
    --weight full_sft --hidden_size 768
\end{verbatim}

\subsection{Verification Harness}
\label{sec:appendix:verify}

Before reporting any results, run the 6-gate verification harness:

\begin{verbatim}
# For CUDA / H100
bash paper/experiments/verify_all.sh cuda

# For ROCm / MI350
bash paper/experiments/verify_all.sh rocm
\end{verbatim}

Expected terminal output on a clean run:

\begin{verbatim}
[PASS] Gate 1: All banners valid
[PASS] Gate 2: Convergence parity within 2%
[PASS] Gate 3: MoE aux_loss values are sane
[PASS] Gate 4: Checkpoint round-trip passed
[PASS] Gate 5: All pytest tests passed
[PASS] Gate 6: Greedy parity passed
==================================
Verification summary for <vendor>
  PASSED: 6
  FAILED: 0
==================================
All gates passed. Proceed with experiments.
\end{verbatim}

\subsection{Expected Device Banner}
\label{sec:appendix:banner}

Every training run emits a device banner at initialization via
\texttt{print\_device\_info()} (\texttt{trainer/device\_utils.py}).
The banner format is:

\begin{verbatim}
[device] NVIDIA H100 SXM5 80GB | CUDA 12.x | torch 2.10.0+cu12x
         rank=0/8 | bf16=True | dist_backend=nccl
         mem_total=81920 MiB
\end{verbatim}

For AMD:

\begin{verbatim}
[device] AMD Instinct MI350 | ROCm 7.1 | torch 2.10.0+rocm7.1
         rank=0/8 | bf16=True | dist_backend=nccl
         mem_total=<PH> MiB
\end{verbatim}

Gate~1 of the verification harness greps all Block~A logs for the
\texttt{[device]}, \texttt{bf16=True}, and \texttt{dist\_backend=nccl}
strings to confirm that each log corresponds to the intended platform.

\subsection{MFU Calculation Details}
\label{sec:appendix:mfu}

Model FLOPs Utilization is computed as:
\[
  \text{MFU} = \frac{6 \times P \times T_{\text{step}} / t_{\text{step}}}%
                    {F_{\text{peak}}}
\]
where $P$ is the number of active parameters (equal to total parameters for
dense models; equal to $P_{\text{active}}$ for MoE), $T_{\text{step}}$ is
the number of tokens per step (batch size $\times$ sequence length $\times$
number of GPUs), $t_{\text{step}}$ is the wall-clock step time in seconds,
and $F_{\text{peak}}$ is the vendor-advertised peak throughput in FLOPs/s
(H100: 989\,TFLOPs/s bf16; MI350: \PH{}\,TFLOPs/s bf16).

For the dense MiniMind-768 model:
$P_{\text{dense}} \approx 64.7\times10^6$ parameters.

For the MoE MiniMind-768 model with 4 experts, 1 active, and 1 shared:
$P_{\text{active}} = P_{\text{base}} + P_{\text{shared}} + P_{\text{1 routed expert}}
\approx \PH{}\times10^6$ parameters.

The \texttt{paper/experiments/compute\_mfu.py} script automates this
calculation from training logs and \texttt{MiniMindConfig} parameters.

\subsection{Metric Collection}
\label{sec:appendix:metrics}

After all experiment logs are collected:

\begin{verbatim}
# Parse all logs, compute MFU, aggregate metrics
python paper/experiments/collect_metrics.py \
    --log_dir paper/logs \
    --out_dir paper/tables

# Outputs:
#   paper/tables/metrics.json  — full structured results
#   paper/tables/metrics.csv   — flat CSV for spreadsheet import
\end{verbatim}

The JSON schema is:

\begin{verbatim}
{
  "block_a": {
    "<vendor>_<tag>": {
      "tokens_per_s": float,
      "mfu_pct": float,
      "peak_mem_gb": float,
      "epoch_min": float,
      "final_loss": float
    }, ...
  },
  "block_e": {
    "<vendor>_<engine>_<concurrency>": {
      "ttft_p50_ms": float,
      "ttft_p95_ms": float,
      "itl_p50_ms": float,
      "tokens_per_s": float,
      "requests_per_s": float
    }, ...
  }
}
\end{verbatim}
